% 本文件由 LaTeX 排版 Agent 根据有效 translation.json 自动生成。
% author=Apocrypha; work_id=0847; title=雅各伯原始福音

\chapter{雅各伯原始福音}

% structure: p0001 source=h1 target=omitted reason=page-title-already-rendered-by-parent

\Emph{天主圣母玛利亚及极荣耀的耶稣基督之母的诞生。}\par

1. 在以色列十二支派的记录中，有一位极其富有的约阿基姆。他献双倍的祭品，说：愿我多余的财物归于全体百姓，愿我的赎罪祭献于上主，以平息祂的义怒。因为上主的大日临近，以色列子民都在献祭。鲁宾站在他对面说：你不配先献祭，因为你在以色列中没有留下后裔。约阿基姆极其忧伤，前往百姓十二支派的登记册处，说：我要查看以色列十二支派的登记册，是否唯独我在以色列中没有后裔。他查考后，发现所有义人都在以色列中留下了后裔。他想起先祖亚巴郎，天主在最后的日子赐给他一个儿子依撒格。约阿基姆极其忧伤，没有与妻子见面，却退到旷野，在那里搭起帐篷，禁食四十昼夜，心里说：我绝不下山吃喝，直到上主我的天主垂顾我；祈祷将成为我的饮食。\par

2. 他的妻子亚纳在双重的哀悼中哀悼，在双重的悲叹中悲叹，说：我要哀哭我的寡居，我要哀哭我的无子。上主的大日临近了；她的使女犹滴说：你还要贬抑你的灵魂多久？看，上主的大日临近了，你哀悼是不合法的。但你拿这条头带，是那位妇女做的给我的；因为我不该佩戴它，因我是使女，而它有君王的样子。亚纳说：离开我！因为我没有做这样的事，上主使我极其卑微。我恐怕是某个恶人给了你，你使我有份于你的罪。犹滴说：我为何要诅咒你？既然上主封闭了你的子宫，不使你在以色列中结出果实。亚纳极其忧伤，脱下丧服，洗净头，穿上婚服，约在第九时辰下到花园散步。她看见一棵月桂树，坐在下面，向上主祈祷说：我们祖先的天主，求祢降福我，垂听我的祈祷，就如祢降福了撒辣的子宫，赐给她儿子依撒格一样。\par

3. 她仰望天空，看见月桂树上有一个麻雀巢，\BibleRef{多 2:10}内心发出悲叹，说：哀哉！谁生了我？哪个子宫孕育了我？因为我成了以色列子民面前的咒骂，我受责斥，他们讥笑着把我赶出上主的圣殿。哀哉！我比作什么？我不像天上的飞鸟，因为天上的飞鸟在祢面前都能生育，上主啊。哀哉！我比作什么？我不像地上的走兽，因为地上的走兽在祢面前都能生育，上主啊。哀哉！我比作什么？我不像这些水，因为连这些水在祢面前都能生育，上主啊。哀哉！我比作什么？我不像这大地，因为连大地都按时结果实，并祝福祢，上主啊。\par

4. 看，有一位上主的天使站在旁边，说：亚纳，亚纳，上主垂听了你的祈祷，你将怀孕，将生育；你的后裔将传遍全世界。亚纳说：上主我的天主永在，我若生男或生女，我必将他作为礼物献给上主我的天主；他一生都要在圣事中事奉祂。\BibleRef{撒上 1:11}看，有两位天使来对她说：看，你的丈夫约阿基姆正带着羊群回来。因为有一位上主的天使下到他那里，说：约阿基姆，约阿基姆，上主天主垂听了你的祈祷。你从这里下去；看，你的妻子亚纳将要怀孕。约阿基姆下去，叫来他的牧人，说：给我带十只无瑕无疵的母羔羊，它们要归于上主我的天主；再带十二只柔嫩的小牛，它们要归于司祭和长老；还有一百只山羊，给全体百姓。看，约阿基姆带着羊群来了；亚纳站在门口，看见约阿基姆来了，就跑上去搂住他的脖子，说：现在我知道上主天主大大降福了我；因为看，寡妇不再寡居，我本无子，将要怀孕。约阿基姆那天第一天在家中休息。\par

5. 次日，他献上祭品，心里说：若上主天主施恩于我，司祭额上的金牌必向我显明。约阿基姆献上祭品，当他走上主的祭台时，他仔细观察司祭的金牌，没有在自己身上看见罪。约阿基姆说：现在我知道上主施恩于我，赦免了我所有的罪。他成义后从上主的圣殿下来，回到自己家中。她的月份满了，在第九个月亚纳生下一女。她对接生婆说：我生了什么？接生婆说：一个女孩。亚纳说：我的灵魂今日被高举。她将她放下。日子满了，亚纳行了取洁礼，给婴儿哺乳，给她起名叫玛利亚。\par

那孩子一天天强壮起来；当她六个月大时，母亲把她放在地上，试试她能否站立，她走了七步，便扑进母亲的怀里；母亲把她抱起，说：\Quote{我指著我的主、天主的永生起誓，你必不得在这地上行走，直到我带妳进入主的圣殿。}她在自己的卧室内设置了一间圣所，不许任何凡俗或不洁之物经过那里。她又召来希伯来人的无玷女儿们，她们却把她引入歧途。当孩子满一岁时，约阿金设了大宴，邀请了司祭、经师、长老和全体以色列民众。约阿金把孩子带到司祭面前；他们祝福她说：\Quote{我们祖先的天主，请祝福这个孩子，赐给她一个永垂不朽的名字，世世代代被传颂。}全体民众都说：\Quote{阿们！阿们！}他又把她带到司祭长面前；他们祝福她说：\Quote{至高的天主，请垂顾这个孩子，以至极的祝福祝福她，直到永远。}母亲把她抱起，带回自己卧室的圣所中，给她喂奶。亚纳向主天主咏唱说：\Quote{我要向我的主、天主歌唱，因为他垂顾了我，除去了我敌人的羞辱；主赐给我义德的果实，独一无二，在他面前充满恩赐。谁能告诉鲁宾的子孙，亚纳在哺乳？听啊，听啊，以色列十二支派，亚纳在哺乳。}她将孩子安放在圣所的卧室内，出去侍奉他们。晚宴结束后，他们欢欣地下去，光荣以色列的天主。\par

孩子的月份逐渐增加。孩子两岁时，约阿金说：\Quote{让我们带她上主的圣殿，好还我们所许的愿，免得主降罚于我们，我们的祭献不被收纳。}亚纳说：\Quote{再等第三年吧，免得孩子寻找父母。}约阿金说：\Quote{那就等吧。}孩子三岁时，约阿金说：\Quote{召来希伯来人的无玷女儿们，让她们每人拿著一盏灯，燃著灯站立，免得孩子回头，心思被引离主的圣殿。}他们就这样行，直到上到主的圣殿。司祭接待了她，亲吻她，祝福她说：\Quote{主在万代中高举了你的名字。在末后的日子，主将藉你向以色列子孙显示他的救赎。}他把她放在祭台的第三级上，主天主将恩宠赐予她；她以脚舞蹈，以色列全家都喜爱她。\par

她的父母惊奇地下去了，赞美主天主，因为孩子没有回头。玛利亚在主的圣殿中，仿佛一只鸽子居住在那里，她从天主使者手中领取食物。当她十二岁时，司祭们召开了会议，说：\Quote{看，玛利亚在主的圣殿中已满十二岁。我们该拿她怎么办呢？免得她玷污主的圣所。}他们对大司祭说：\Quote{你站在主的祭台旁，进去为她祈祷；凡主所指示你的，我们必遵行。}大司祭进去，披著带有十二个铃铛的礼服，进入至圣所，为她祈祷。看，主的天使站在他身旁，对他说：\Quote{匝加利亚，匝加利亚，出去召集民众中的鳏夫，让他们每人拿著自己的杖；主向谁显示记号，谁就娶她为妻。}传令官出去，遍巡犹太全境，主的号角吹响，众人都跑来。\par

若瑟丢下斧头，出来迎接他们；他们聚集后，带著自己的杖去见大司祭。大司祭接过所有人的杖，进入圣殿祈祷；祈祷完毕，他拿出杖来，还给他们：杖上毫无记号，若瑟最后拿自己的杖；看，一只鸽子从杖中飞出，落在若瑟的头上。大司祭对若瑟说：\Quote{你已被抽中，要保管主的童女。}但若瑟拒绝说：\Quote{我有儿女，而且我已年老，她却是年轻女子。我恐怕会成为以色列子孙的笑柄。}大司祭对若瑟说：\Quote{你要敬畏你的主、天主，要记住主对达堂、阿彼兰和科辣黑所行的事；\BibleRef{户 16:31}大地怎样裂开，因他们的反抗而吞没了他们。如今，若瑟，你要害怕，免得同样的事发生在你家中。}若瑟害怕了，便把她置于自己保管之下。若瑟对玛利亚说：\Quote{看，我从主的圣殿中接受了妳；如今我把妳安置在我家中，然后我要去建造我的建筑物，再回到妳这里。主必保护妳。}\par

当时司祭们召开了会议，说：\Quote{让我们为主的圣殿做一幅帷幔。}司祭说：\Quote{把达味家族中无玷的童贞女召到我这里来。}差役出去寻找，找到了七位童贞女。司祭想起了孩子玛利亚，她是达味家族的，在天主面前是无玷的。差役出去把她带来。他们带她们进入主的圣殿。司祭说：\Quote{为我抽签决定谁纺金线、白线、细麻、丝绸、蓝线、朱红线、真紫色线。}\BibleRef{出 25:4}真紫色线和朱红线抽中了玛利亚，她便拿了这些，回到自己家中。那时匝加利亚成了哑巴，撒慕尔代替他的位置，直到匝加利亚开口说话。玛利亚拿了朱红线，纺了起来。\par

她拿起水罐，出去打水。看，有一个声音说：\BibleQuote{蒙受恩宠的，祝你平安！主与你同在！你在妇女中应受赞颂！}\BibleRef{路 1:28}她左右观看，要看这声音从何处来。她战战兢兢地回到家中，放下水罐；拿起紫线，坐在座位上，将它抽出来。看，主的天使站在她面前，说：\Quote{玛利亚，不要怕！因为你在万主之主面前蒙受了恩宠，你将按祂的话怀孕。}她听见了，心中思忖说：\Quote{我岂能由主，生活的天主怀孕？我岂能像一般妇女那样生产？}主的天使说：\Quote{玛利亚，不是这样；因为主的能力要荫庇你，因此那要诞生的圣者，将称为至高者的儿子。你要给祂起名叫耶稣，因为祂要把自己的民族从他们的罪恶中拯救出来。}玛利亚说：\Quote{看，上主的婢女在祂面前，愿照你的话成就于我吧！}\par

她完成了紫线和朱红线，拿去交给司铎。司铎祝福了她，说：\Quote{玛利亚，上主天主已使你的名显大，你在大地的一切世代中将受赞颂。}玛利亚极其喜乐，往她亲戚依撒伯尔那里去（\BibleRef{路 1:39}），敲了门。依撒伯尔听见她的声音，便丢下朱红线，跑到门前，开了门；看见玛利亚，就祝福她说：\BibleQuote{我主的母亲到我这里，这事怎能临到我？因为，看，我腹中的胎儿跳跃着祝福了你。}\BibleRef{路 1:34}但玛利亚已忘记了总领天使加俾额尔所说奥秘的事，仰天注视，说：\BibleQuote{主啊！我算什么，竟使大地一切世代都祝福我？}\BibleRef{路 1:48}她与依撒伯尔同住了三个月，日子一天天过去，她的身体渐渐显怀。玛利亚害怕，便回到自己的家中，躲在以色列子民面前。这些奥秘的事发生时，她十六岁。\par

她怀孕六个月了；看，若瑟从建筑工地回来，进了他的家，发现她已经身怀六甲。他便拍打自己的脸，扑倒在地上，穿着苦衣，痛哭流涕，说：\Quote{我怎能仰面见上主我的天主？关于这个少女，我该做什么祈祷？因为我从主的圣殿领受了她，她是童贞女，而我却没有看顾好她。是谁暗算了我？是谁在我家中行了这恶事，污辱了这童贞女？亚当的故事岂不重现在我身上？正如亚当正在歌唱赞美时，蛇来了，找到厄娃独自一人，完全欺骗了她，同样的事也发生在我身上。}若瑟从苦衣上站起来，叫了玛利亚，对她说：\Quote{你这蒙天主眷顾的，你为什么做了这事，忘记了上主你的天主？你为什么贬低你的灵魂？你本是在至圣所中养育的，从天使手中领受食物。}她痛哭流涕，说：\Quote{我是无罪的，没有认识男人。}若瑟对她说：\Quote{那么，你腹中的是从哪里来的？}她说：\Quote{我指着上主我的天主起誓，我不知道它从哪里来。}\par

若瑟非常恐惧，便离开她，思虑该如何处置她。（\BibleRef{玛 1:19}）若瑟说：\Quote{我若隐瞒她的罪，就是与上主的法律对抗；我若将她暴露在以色列子民面前，又怕她腹中的是来自天使，我将会把无辜的血交给死亡。那么我该怎样对待她？我暗暗地休了她吧。}夜幕降临时，看，主的天使在梦中显现给他，说：\BibleQuote{不要怕这少女，因为她腹中的是出于圣神；她将生一个儿子，你要给祂起名叫耶稣，因为祂要把自己的民族从他们的罪恶中拯救出来。}\BibleRef{玛 1:20}若瑟从睡眠中醒来，光荣了以色列的天主，因祂赐给了他这恩宠，他便留住了她。\par

经师阿纳尼雅来到他那里，说：\Quote{你为什么没有出现在我们的集会中？}若瑟对他说：\Quote{因为我旅途劳顿，第一天休息了。}他转过身，看见玛利亚有孕在身。他便跑到司铎那里，对他说：\Quote{你所担保的那个若瑟，犯了重罪。}司铎说：\Quote{怎么回事？}他说：\Quote{他污辱了他从主的圣殿所领受的童贞女，偷偷娶了她，并没有向以色列子民揭露。}司铎回答说：\Quote{若瑟做了这事？}经师阿纳尼雅说：\Quote{派差役去，你会找到那怀孕的童贞女。}差役去了，果然如他所说；他们便把她和若瑟一同带到法庭。司铎说：\Quote{玛利亚，你为什么做了这事？你为什么贬低你的灵魂，忘记了上主你的天主？你本是在至圣所中养育，从天使手中领受食物，听见赞歌，在祂面前舞蹈，你为什么做了这事？}她痛哭流涕，说：\Quote{我指着上主我的天主起誓，我在祂面前是纯洁的，没有认识男人。}司铎对若瑟说：\Quote{你为什么做了这事？}若瑟说：\Quote{我指着上主起誓，我对她是纯洁的。}司铎便说：\Quote{不可作假见证，要讲实话。你偷偷娶了她，没有向以色列子民揭露，也没有在强有力的手下低头，好使你的后裔蒙祝福。}若瑟默不作声。\par

16. 司祭说：\Quote{把你从主的圣殿接来的童贞女交出来。}若瑟便泪流满面。司祭说：\Quote{我要给你喝主的试炼之水，主要在你的眼中显明你的罪。}司祭取了水，给若瑟喝，并打发他往山地去；他安然无恙地回来了。他也给玛利亚喝，并打发她往山地去；她也安然无恙地回来了。众人都惊奇，因为罪没有在他们身上显现。司祭说：\Quote{如果主天主没有显明你们的罪，我也不审判你们。}遂打发他们走了。若瑟带着玛利亚，欢欢喜喜地回到自己的家，光荣赞美以色列的天主。\par

17. 奥古斯都皇帝下令，凡在犹太的白冷的人都应登记\BibleRef{路 2:1}。若瑟说：\Quote{我要登记我的儿子们，但这少女我该怎么办？我怎么登记她？作为我的妻子？我羞于这样做。那么作为我的女儿？但以色列所有的儿子都知道她不是我的女儿。主的日子自会照主的旨意成就。}他给驴备上鞍，让她骑上；他的儿子牵着驴，若瑟跟在后面。当走了约三哩时，若瑟转过身来，见她面带愁容；他心中暗想：\Quote{大概是她腹中的胎儿使她难受。}若瑟又转过身来，见她面带笑容。他对她说：\Quote{玛利亚，为什么我看见你脸上时而欢笑，时而忧愁？}玛利亚对若瑟说：\Quote{因为我用眼睛看见两个民族：一个哭泣哀号，另一个欢欣喜乐。}他们走到路中间，玛利亚对他说：\Quote{把我从驴上扶下来，因为我腹中的胎儿急着要出来。}他便把她从驴上扶下来，对她说：\Quote{我带你去哪里？在哪里遮掩你的羞耻？因为这里是旷野。}\par

18. 他在那里找到一个山洞，带她进去；把两个儿子留在她身边，自己出去在白冷地区寻找一位产婆。\par

我若瑟走着，却又像没有走；我抬头望天，见天空惊愕；我仰望天穹，见它静止不动，空中的飞鸟也静止了。我低头看地，见一个槽子放着，工人们斜躺着：他们的手在槽里。吃饭的人没有吃，起身的人没有起身，送食物到口边的人没有送；所有人的脸都向上仰望。我看见羊在行走，羊却停住了；牧羊人举手要打羊，手却举在空中。我看河水的流势，见山羊羔的嘴停在水上不喝，一切事物瞬间都脱离了它们的轨道。\par

19. 我看见一个女人从山地下来，她对我说：\Quote{人啊，你往哪里去？}我说：\Quote{我在找一位希伯来产婆。}她回答我说：\Quote{你是以色列人吗？}我对她说：\Quote{是。}她说：\Quote{那在山洞里生产的女子是谁？}我说：\Quote{是一位许配给我的女子。}她对我说：\Quote{她不是你的妻子吗？}我对她说：\Quote{她是玛利亚，在主的圣殿中长大，我抽签得了她为妻。但她并不是我的妻子，却因圣神而怀孕。}\par

产婆对他说：\Quote{这是真的吗？}若瑟对她说：\Quote{你来看看。}产婆便跟他去了。他们站在山洞所在的地方，看哪，一朵发光的云彩笼罩了山洞。产婆说：\Quote{今天我灵被高举，因为我亲眼看见了奇事——救恩已为以色列诞生了。}云彩立刻从山洞消失，山洞里大放光明，眼睛不能承受。稍后，光明渐渐减弱，直到婴孩出现，并走去从祂母亲玛利亚那里吸吮乳汁。产婆呼喊说：\Quote{今天对我来说是伟大的一天，因为我看见了这奇异的景象。}产婆从山洞出来，撒罗默遇见她。她对她说：\Quote{撒罗默，撒罗默，我有件奇事要告诉你：一个童贞女生产了——这是她的本性所不容的。}撒罗默就说：\Quote{我指着主我的天主起誓，除非我用手指探入并检查，我绝不相信童贞女会生产。}\par

20. 产婆进去，对玛利亚说：\Quote{显露你自己吧；因为关于你起了不小的争论。}撒罗默伸进手指，喊叫说：\Quote{我有祸了！为我的不义和不信，因为我试探了永生的天主；看哪，我的手像被火烧着一样掉下来了。}她跪在主前，说：\Quote{我祖先的天主啊，求你记念我是亚巴郎、依撒格、雅各伯的后裔；不要使我在以色列子民面前蒙羞，但求你把我恢复给穷人；因为你知道，主啊，我奉你的名尽了职，也从你手中得了赏报。}看哪，主的一位天使站在她旁边，对她说：\Quote{撒罗默，撒罗默，主垂听了你。把你的手放到婴孩上，抱着祂，你必得平安和喜乐。}撒罗默便去抱起祂，说：\Quote{我要朝拜祂，因为一位伟大的君王已为以色列诞生了。}看哪，撒罗默立刻痊愈了，她便从山洞出来，成了义人。看哪，有声音说：\Quote{撒罗默，撒罗默，不要把你所见的奇事告诉人，直到孩子进入耶路撒冷。}\par

21. 看，若瑟正准备往犹太去。在犹太的白冷发生了一场大骚动，因为贤士们来了，说：\Quote{那生下来作犹太人之王的在哪里？我们在东方看见了他的星，特来朝拜他。}黑落德听见了，就十分惊慌，便派差役到贤士那里去。他又召来司祭，查问他们说：\Quote{论到基督，是怎么写的？祂要在哪里出生？}他们回答说：\Quote{在犹太的白冷，因为经上这样记载。}他就打发他们走了。他又查问贤士们，对他们说：\Quote{你们看见了什么关于那已出生的王的异兆？}贤士们说：\Quote{我们看见一颗极大的星照耀在这些星之间，遮蔽了它们的光，以致众星都不显现；这样我们就知道一位君王已在以色列出生，特来朝拜他。}黑落德说：\Quote{你们去寻访那孩子，若找到了，就报告给我，好叫我也去朝拜他。}贤士们便出发了。看，他们在东方所看见的那颗星，在他们前头走，直到来到山洞上方，就停在洞顶。贤士们看见婴孩与祂的母亲玛利亚在一起，就从袋中取出黄金、乳香和没药。他们由于在梦中被天使警告不要回犹太去，就从另一条路返回了自己的家乡。\par

22. 黑落德知道自己被贤士们愚弄了，就大怒，派遣凶手，对他们说：\Quote{凡两岁以下的婴孩，都杀掉。}玛利亚听见孩子正被杀害，就害怕，抱起婴孩，裹好了祂，把祂放在一个牛棚里。依撒伯尔听见他们正在搜捕若翰，就带着他上了山地，一直寻找可以藏匿他的地方，却没有可藏之处。依撒伯尔大声叹息说：\Quote{天主的山啊，收纳母亲和孩子吧！}立时那山裂开，收纳了她。有一道光在他们周围照耀，因为上主的一位天使与他们同在，守护着他们。\par

23. 黑落德搜捕若翰，就派差役到匝加利亚那里，说：\Quote{你把你的儿子藏在哪里了？}匝加利亚回答他们说：\Quote{我是天主在圣事上的仆人，常在上主的圣殿里侍立；我不知道我的儿子在哪里。}差役离去，将这一切事报告给黑落德。黑落德大怒，说：\Quote{他的儿子注定要作以色列的王。}于是他又派人去对匝加利亚说：\Quote{你要说实话，你的儿子在哪里？你知道你的性命在我手中。}匝加利亚说：\Quote{我是天主的殉道者，你即使流我的血；因为主要收纳我的灵魂，你在上主圣殿的门廊流了无辜者的血。}匝加利亚约在破晓时被杀害。以色列的子孙不知道他被杀害了。\par

24. 在问候的时辰，司祭们离去，匝加利亚没有照常出来以祝福迎接他们。司祭们站着等候匝加利亚在祈祷时问候他，并光荣至高者。他迟迟不来，众人都害怕了。其中一人壮胆走进去，看见祭台旁边有凝血；他听见有声音说：\Quote{匝加利亚已被杀害，他的血不得擦去，直到为他复仇的人来到。}他听见这话，就害怕，出来告诉了司祭们。他们壮胆进去，看见了所发生的事；圣殿的格子结构发出哀鸣之声，他们从顶到底撕裂了衣服。他们找不到他的尸体，却发现他的血变成了石头。他们害怕了，出去报告民众说匝加利亚被杀害了。民众的众支派听见了，就哀悼他三天三夜。三天后，司祭们商议该立谁在他位上；掣签选出了西默盎。因为圣神曾启示他，在未见到基督以肉身显现之前，他决不会见死亡。\par

而我——雅各伯，在耶路撒冷写成了这篇历史；当黑落德死时，发生了骚动，我就退到旷野，直到耶路撒冷的骚动平息，光荣了主天主，祂赐给了我恩赐和智慧来写这篇历史。愿恩宠归于敬畏我们主耶稣基督的人；愿光荣归于祂，至于无穷之世。阿们。\par

\SourceRecord{Translated by Alexander Walker.；Ante-Nicene Fathers；Vol. 8.；Edited by Alexander Roberts, James Donaldson, and A. Cleveland Coxe.；Buffalo, NY: Christian Literature Publishing Co.,；1886.；<http://www.newadvent.org/fathers/0847.htm>.}
